\section{Selection, Provenance, and Reading Depth}\label{app:selection}
\subsection{Two evidence layers}
The contextual library supports a narrative account; the documentary audit supplies the comparative cases in Section~\ref{sec:analysis}. Their populations and reading depths must not be conflated. The audit includes AlignScore, FActScore, SelfCheckGPT, semantic entropy, Lookback Lens, FactAlign, Self-RAG, CAD, DoLa, TruthX, CoVe, CRITIC, and SEAL. These were all the named studies in a pre-existing resource table, chosen for mechanism diversity. The sample was fixed before the new extraction round on 5 October 2026; this is a recorded protocol, not an externally preregistered study.

We recorded primary publication URLs, PDF digests, inspected sections, and claim-specific locators. Extraction fields cover reference frame, scoring unit, evidence/access, information control, feedback--evaluation separation, resources, and joint validation. Information controls are described rather than collapsed into an adequate/inadequate score: fact quantity, task quality, required-aspect coverage, and answered-question coverage measure different things. Different prompts or model names do not prove statistically independent errors. A qualitative access requirement is not measured latency. ``Not located'' is bounded by the inspected source passages and is not proof of absence. The supplementary CSV preserves caveats for every case. Source checks and revisions were AI-assisted; no independent human annotation or inter-annotator agreement is claimed.

\subsection{Contextual discovery and targeted update}
The original contextual metadata snapshot is 15 September 2026. Discovery queried official ACL Anthology collections across ACL, EMNLP, NAACL, EACL, Findings, and TACL in 2022--2026. Twenty-five collections were available and five requested combinations unavailable. A broad title filter returned 1,766 records; three targeted additions and seven external studies brought the candidate pool to 1,776. The 87-study working library was selected for relevance, mechanism diversity, benchmark coverage, and counterevidence. Automatic title triage excluded 401 candidates; 1,288 remained unassessed reserve records. All selected titles and abstracts were screened; initial preparation inspected passages in fifteen papers, followed by claim-specific checks. These are workflow counts, not a full-text systematic-review flow.

The targeted update on 5 October 2026 checked closely related surveys and evaluation work against primary publication records. Its query and inclusion log is supplied separately; it does not purport to exhaust all intervening publications. Canonical published versions are preferred. The audit sample remains fixed, so update records and the broader bibliography do not increase its denominator. No model experiments, independent benchmark replications, or pooled effect estimates were performed.

\subsection{Visual provenance and coding examples}
The layouts of Figures~\ref{fig:examples}, \ref{fig:evaluation}, and~\ref{fig:mitigation} and Table~\ref{tab:hallucination-benchmarks} follow \citet{liu-etal-2024-lvlm-design}; Figure~\ref{fig:structure} follows \citet{liu-etal-2025-logical-design}. Their content addresses text LLMs. Figure~\ref{fig:examples} uses synthetic interactions. Figure~\ref{fig:mitigation} summarizes mechanisms proposed in the literature, not an experimentally identified causal decomposition. The figures provide a representative conceptual map; the source-linked audit adds protocol details that cannot be inferred from this map.

Detection branches are assigned by decision evidence, mitigation branches by the operation that changes generation. Thus Lookback Lens appears as attention-based detection and candidate-selection mitigation. RHIO combines training on induced negatives with contrastive decoding. Hybrid methods need not belong to exclusive categories. Source faithfulness is recorded separately from world factuality, and target-generator access separately from proxy or verifier access.

\subsection{Compatibility audit records}\label{app:audit}
The accompanying CSV contains one record per fixed case, with primary-source URLs, PDF digests, inspected sections, and locators for each substantive comparison. The records describe concrete reported controls rather than assign a universal quality score. They distinguish measured costs from access requirements, human labels from model judgments, and component analyses from tests of the complete detector--intervention system. Table~\ref{tab:method-comparison} highlights contrasts used in the main argument. Other selected studies remain in the record even when they do not support a simple positive or negative code. Tables~\ref{tab:audit-1}--\ref{tab:audit-3} summarize reported information controls and their boundaries.
\begin{table*}[tp]
\centering\small\setlength{\tabcolsep}{5pt}\renewcommand{\arraystretch}{1.12}
\begin{tabularx}{\textwidth}{@{}>{\raggedright\arraybackslash}p{.18\textwidth}>{\raggedright\arraybackslash}p{.42\textwidth}>{\raggedright\arraybackslash}X@{}}\toprule
Study & Reported information control & Interpretation boundary \\ \midrule
AlignScore \citep{zha-etal-2023-alignscore} & Fixed-output metric evaluation. Explicitly separates omission/content-selection errors from factual-consistency errors in PolyTope coding. A required-aspect-completeness or selective-answering outcome measure was not located in the inspected evaluation sections; omission exclusion is not such a measure. & Factuality and information selection are deliberately separated; detector performance does not alone establish usefulness after applying an output intervention. The altered PolyTope error definition and dataset overlap exclusions must be retained in comparisons. \\
FActScore \citep{min-etal-2023-factscore} & Reports responding fraction and fact quantity per response; explicitly warns that precision does not measure recall. The related editing experiment measures token preservation and agreement with human edits. Fact quantity and edit similarity are not required-aspect completeness. & Unsupported includes source absence as well as contradiction. Responding fraction and fact counts partly expose selective/short responses but cannot establish complete answers. Oracle-label editor results must not be attributed to the automatic estimator. \\
SelfCheckGPT \citep{manakul-etal-2023-selfcheckgpt} & Fixed-output detection/ranking on generated biographies. Descriptive token/sentence statistics are reported, but matched answered-question retention or required-aspect-completeness evaluation was not located in inspected \S{}\S{}6--7 and Appendix D. & Sample consistency can miss stable incorrect beliefs; the observed task is a selected GPT-3 biography set. Sentence PR-AUC and passage correlation do not measure correctness/completeness after editing or abstention. \\
Semantic entropy \citep{extsemanticentropy2024} & Explicit answered-question retention through rejection-accuracy curves and AURAC; FactualBio also evaluates retained fact predictions. This controls a retained-unit fraction, not required-aspect completeness of long answers. Biography question generation attempts to cover claim aspects but has acknowledged misses. & QA labels operationalize reference-answer agreement, not independently established world truth; SI explicitly preserves that target when a reference is wrong. FactualBio has only 21 biographies/150 claims. Retained-unit coverage differs from semantic completeness; systematic low-entropy errors remain outside the targeted confabulation mechanism. \\
Lookback Lens \citep{chuang-etal-2024-lookback} & Reports a general task-quality proxy (original MT-Bench score) alongside contextual-faithfulness score and NQ exact match. These do not establish matched required-aspect coverage or answer-completeness preservation for XSum. & The quality safeguard is task-specific and is not an equivalence test or matched-completeness study. Shared GPT-4o labeling/judging can preserve common errors despite human audit. Resource figures are approximate study-specific observations; transfer needs an auxiliary fitted head mapping. \\
\bottomrule\end{tabularx}
\caption{Documentary audit of reported controls (1/3). Descriptions retain original-study targets and unresolved reporting; precise passage locators and source digests are in the accompanying evidence records.}
\label{tab:audit-1}\end{table*}

\begin{table*}[tp]
\centering\small\setlength{\tabcolsep}{5pt}\renewcommand{\arraystretch}{1.12}
\begin{tabularx}{\textwidth}{@{}>{\raggedright\arraybackslash}p{.18\textwidth}>{\raggedright\arraybackslash}p{.42\textwidth}>{\raggedright\arraybackslash}X@{}}\toprule
Study & Reported information control & Interpretation boundary \\ \midrule
FactAlign \citep{huang-chen-2024-factalign} & Fact counts, factual precision and f1@K; literal Section 3.1 recall term is min(1, number of generated atomic statements/K), not a reference-aspect recall. MT-Bench provides helpfulness proxy. No matched answered-question coverage comparison located in inspected sections. & Corpus support is not universal truth. Generated fact quantity is not required-information completeness. For loss ablation compare the one-iteration full model with its ablation, not the three-iteration headline result. \\
Self-RAG \citep{extselfrag2024} & ASQA correctness, ROUGE, MAUVE and citation precision/recall; biography factual precision; ASQA maximum 100 generated tokens across baselines to limit length bias. Weight sweep exposes support-versus-fluency trade-off. These are not matched answered-question coverage or required-aspect completeness controls. & S\&P denominator is conditional on the model support/relevance decisions. Reflection agreement with GPT-4 is not independent truth adjudication. Citation recall concerns generated claims, not all information a user requires. \\
CAD \citep{shi-etal-2024-trusting} & Summarization reports ROUGE-L alongside FactKB and BERT-precision; QA reports exact match. Task quality is monitored, but neither reference-aspect completeness nor matched answered-question coverage is established. & NQ-Swap intentionally changes source entities, so success may require contradicting world knowledge. Do not classify its exact-match gain as world-truth improvement or infer a twofold latency from two distributions. \\
DoLa \citep{extdola2024} & TruthfulQA reports Truth, Info, their joint metric, and Reject; GPT-4 separately rates grammaticality/cohesiveness. These show informativeness/refusal trade-offs but are not matched coverage. For LLaMA-33B, Truth falls 62.5 to 56.4 while Info rises 69.0 to 92.4 and joint score rises 31.7 to 49.1. & Cannot add unlearned factual knowledge. Raw truth and informativeness may move in opposite directions. Latency result belongs to the reported hardware/output-length configuration. \\
\bottomrule\end{tabularx}
\caption{Documentary audit of reported controls (2/3). Descriptions retain original-study targets and unresolved reporting; precise passage locators and source digests are in the accompanying evidence records.}
\label{tab:audit-2}\end{table*}

\begin{table*}[tp]
\centering\small\setlength{\tabcolsep}{5pt}\renewcommand{\arraystretch}{1.12}
\begin{tabularx}{\textwidth}{@{}>{\raggedright\arraybackslash}p{.18\textwidth}>{\raggedright\arraybackslash}p{.42\textwidth}>{\raggedright\arraybackslash}X@{}}\toprule
Study & Reported information control & Interpretation boundary \\ \midrule
TruthX \citep{zhang-etal-2024-truthx} & Truth, Info, joint score and no-comment counts; original-model perplexity and cross-task accuracy supplement quality. They monitor refusal/informativeness, not matched answered coverage or required information completeness. & Two-fold in-domain validation and multiple-choice transfer do not establish general free-generation truth calibration. Editing cannot guarantee truth or inject missing knowledge (Limitations p.10). \\
CoVe \citep{dhuliawala-etal-2024-chain} & Correct/incorrect entity counts, multi-span precision/recall/F1, biography fact count with FActScore. Same-generator biography result 55.9 to 71.4 FActScore accompanies 16.6 to 12.3 facts. A separate Llama-2-chat clipping control is not a matched-length CoVe control. & Precision gain does not establish preservation of all correct or required facts. Closed-book checks remain bounded by model knowledge; shorter output is a measured trade-off, not proof of the entire gain being due to deletion. \\
CRITIC \citep{extcritic2024} & QA EM/F1 and qualitative/manual refusal and incorrect-correction categories; math answer correctness; toxicity task reports fluency/diversity proxies. No matched answered-question coverage or required-aspect preservation located. Non-oracle CRITIC and oracle CRITIC* must be kept separate. & Table 6 error-type percentages have a less explicit denominator than a full transition matrix; do not treat 10\% incorrect correction as the population rate. Oracle CRITIC* is not deployable verification; math/toxicity results do not directly measure factual hallucination. \\
SEAL \citep{huang-etal-2025-alleviating} & Biography correct-claim counts and FActScore; LongFact precision/R@48/F1@48; separate refusal experiment reports known-question accuracy versus unknown-question rejection across tau; instruction following is checked. R@48 uses truthful-fact count relative to 48, not required-aspect recall or matched question coverage. & Known/unknown status is a correctness/popularity proxy, not direct access to parametric knowledge. Main result cannot be interpreted as answer abstention. Biography sample count conflicts between Table 8 (138) and Appendix B.2 (183), unresolved. \\
\bottomrule\end{tabularx}
\caption{Documentary audit of reported controls (3/3). Descriptions retain original-study targets and unresolved reporting; precise passage locators and source digests are in the accompanying evidence records.}
\label{tab:audit-3}\end{table*}


\FloatBarrier

\subsection{Comparison with prior surveys}\label{app:prior}
The source-linked prior-survey record covers general NLG hallucination, LLM factuality, broad LLM hallucination, human-evaluation reporting, and detection/mitigation reviews. Its purpose is to delimit the contribution, not establish that a concept is absent from an entire literature. The resulting contribution claim is a cross-stage compatibility synthesis applied to a bounded sample. No first-survey or first-audit claim is made.


\section{Benchmark Units and Protocol Caveats}\label{app:benchmarks}
Table~\ref{tab:benchmark-caveats} supplies the counting and interpretation details behind Table~\ref{tab:hallucination-benchmarks}. A benchmark may support both a description of generated errors and evaluation of a detector, but those uses require different denominators and sampling assumptions. In particular, reporting error frequencies within deliberately selected difficult examples does not estimate a generator's error prevalence on ordinary requests.

\begin{table*}[tp]
\centering\small
\setlength{\tabcolsep}{5pt}\renewcommand{\arraystretch}{1.14}
\begin{tabularx}{\textwidth}{@{}p{.17\textwidth}X@{}}\toprule
Benchmark & Unit, protocol, and limitation \\\midrule
TruthfulQA & The original 2022 benchmark contains 817 questions. Generation measures truthfulness/informativeness; MC1 measures single-true accuracy, and MC2 normalized probability mass assigned to true answers. The \href{https://github.com/sylinrl/TruthfulQA}{official January 2025 update} introduced two-option choices and revised the question set. Report the dataset version and MC protocol separately; answer selection is not detector classification. \\
HaluEval & The 35,000 examples combine task-specific generated examples with a human-annotated general-query portion. Artificial and naturally generated errors should be reported separately rather than treated as one deployment distribution. \\
ANAH & Approximately 12,000 annotated sentences provide error types, evidence, and corrections. Generative versus discriminative annotators describe detector implementations; their task remains prediction of annotation labels. \\
RAGTruth & The 17,790 responses come from 2,965 prompts and six generators, not 17,790 independent documents. Response classification and span localization have different matching rules and must be evaluated separately. \\
ExpertQA & The 2,177 expert-curated questions support factuality and attribution analysis. Correctness and cited-source support are separate judgments; exact evaluation subsets and answer-generation conditions matter. \\
TofuEval & Section 3.3 reports 1,500 initial summaries, removal of 23, and 1,479 retained summaries. Those values are arithmetically inconsistent. We retain the reported 1,479 rather than invent a corrected dataset size. Topic relevance and completeness are separate from factual consistency. \\
HALoGEN & The 10,923 prompts cover all nine domains, including programming. That is the complete-suite count, not a verified text-only subset size. Different domain verifiers do not constitute one universal human truth oracle. \\
FaithBench & The 750 summaries derive from 75 passages and ten models, selected using detector disagreement. This supports difficult-case detector analysis; error frequencies are conditional on selection, not population prevalence. \\
FactBench & The 1,000-prompt snapshot uses user-derived prompts and fine-grained verification. Supported, unsupported, and undecidable units require explicit aggregation rules. A dynamic benchmark needs a fixed version for reproducible comparisons. \\
\bottomrule\end{tabularx}
\caption{Evaluation units, protocols, and limitations of the benchmarks in Table~\ref{tab:hallucination-benchmarks}.}
\label{tab:benchmark-caveats}
\end{table*}

\section{Supplementary Comparative Evidence}\label{app:comparisons}
Tables~\ref{tab:extra-evaluation}--\ref{tab:extra-mitigation} extend the main synthesis with complementary studies. The final column states the inference boundary used in this survey; it is not an assertion that the original authors omitted every proposed test. These studies expand coverage of mechanisms and failure conditions rather than form a ranked list. Descriptive entries are grounded in the recorded primary metadata and selected source passages, with reading depth retained in the repository.
\begin{table*}[tp]
\centering\small\setlength{\tabcolsep}{5pt}\renewcommand{\arraystretch}{1.18}
\begin{tabularx}{\textwidth}{@{}p{.22\textwidth}p{.34\textwidth}X@{}}\toprule
Study & Mechanism or evaluation design & Comparative implication and boundary \\\midrule
TRUE \citep{honovich-etal-2022-true-evaluating} & Standardizes factual-consistency datasets across tasks and evaluates metrics at the example level. & Instance discrimination differs from system-level correlation. Combined existing tasks do not estimate real-world hallucination prevalence. \\
AggreFact analysis \citep{tang-etal-2023-understanding} & Stratifies summarization error annotations by generating model and error type; metric gains vary across strata. & Pooled scores can hide weaknesses on newer generators or rare error types. \\
BUMP \citep{ma-etal-2023-bump} & Uses minimally differing human-written summary pairs with one introduced error. & A controlled error-sensitivity probe complements natural generation data; minimal pairs do not reproduce the full distribution of model errors. \\
ALCE \citep{gao-etal-2023-enabling} & Evaluates answer quality and citation quality in end-to-end retrieval and cited generation. & Citation presence, citation support, and factual correctness are distinct criteria. A plausible citation is not by itself evidence. \\
LongEval \citep{krishna-etal-2023-longeval} & Studies fine-grained human judgments and partial annotation for long summaries. & Fine-grained annotation requires agreement and coverage reporting; partial annotation does not exhaustively verify all claims. \\
FACTOR \citep{muhlgay-etal-2024-generating} & Transforms factual corpora into true versus false continuation choices. & Controlled factual choice evaluates a different task from open-ended generation and error localization. \\
ReEval \citep{yu-etal-2024-reeval} & Perturbs evidence to create dynamic tests of evidence use under retrieval. & Static accuracy may reflect memorized knowledge rather than source adherence; counterfactual evidence requires an explicit faithfulness target. \\
Core \citep{jiang-etal-2025-core} & Filters redundant or uninformative subclaims before factual precision aggregation. & Trivial repetitions can inflate precision; informativeness and coverage are complementary criteria. \\
TRIVIA+ \citep{chen-etal-2026-rethinking-evaluation} & Adds long-context RAG detection data and controlled sample-dependent label noise. & Long-context and noisy-label tests complement evaluation on clean labels. \\
RGB \citep{extrgb2024} & Separates noise robustness, negative rejection, information integration, and counterfactual robustness in English/Chinese RAG. & Retrieval conditions define distinct RAG abilities, not one interchangeable hallucination score. \\
\bottomrule\end{tabularx}
\caption{Additional benchmarks and evaluation protocols for hallucination in LLMs.}
\label{tab:extra-evaluation}
\end{table*}

\begin{table*}[tp]
\centering\small\setlength{\tabcolsep}{5pt}\renewcommand{\arraystretch}{1.18}
\begin{tabularx}{\textwidth}{@{}p{.22\textwidth}p{.34\textwidth}X@{}}\toprule
Study & Mechanism or evaluation design & Comparative implication and boundary \\\midrule
WeCheck \citep{wu-etal-2023-wecheck} & Aggregates weak labels over actual generated samples, then learns a noise-aware classifier. & Supervision quality and error distribution distinguish a learned consistency model from zero-shot entailment. \\
FactKB \citep{feng-etal-2023-factkb} & Uses knowledge-base factual pretraining with entity/relation objectives for factuality evaluation. & Structured supervision targets relation errors; broader transfer requires separate evidence. \\
TrueTeacher \citep{gekhman-etal-2023-trueteacher} & Uses LLM-labeled model-generated summaries to train a smaller consistency evaluator. & Distillation trades annotation cost for teacher-label dependence; generated errors differ from artificial perturbations. \\
FENICE \citep{scire-etal-2024-fenice} & Extracts atomic claims and aligns them to a source using natural language inference (NLI). & Local alignment improves interpretability, while claim extraction and entailment both remain possible error sources. \\
FactReasoner \citep{marinescu-etal-2025-factreasoner} & Models evidence--claim logical relationships probabilistically after decomposition/retrieval. & Joint reasoning differs from independently verifying claims; its result still depends on evidence retrieval and relation estimates. \\
Future context \citep{lee-etal-2026-enhancing} & Samples continuations and uses them as clues for black-box hallucination detection. & Generated continuations extend consistency checking without supplying independent external verification. \\
FaithLens \citep{si-etal-2026-faithlens} & Trains detection and explanations using filtered synthetic data, supervised fine-tuning (SFT), and rule-based reinforcement learning. & Explanation quality and decision accuracy require separate validation; fluency alone does not establish explanation faithfulness. \\
QASemConsistency \citep{cattan-etal-2026-localizing} & Decomposes claims into predicate--argument QA pairs for fine-grained support localization. & Semantic error localization supplies a finer target than response-level binary classification. \\
\bottomrule\end{tabularx}
\caption{Additional hallucination detection and explanation methods for LLMs.}
\label{tab:extra-detection}
\end{table*}

\begin{table*}[tp]
\centering\small\setlength{\tabcolsep}{5pt}\renewcommand{\arraystretch}{1.18}
\begin{tabularx}{\textwidth}{@{}p{.22\textwidth}p{.34\textwidth}X@{}}\toprule
Study & Mechanism or evaluation design & Comparative implication and boundary \\\midrule
Factual confidence \citep{mahaut-etal-2024-factual} & Compares confidence estimators and checks stability under meaning-preserving input changes. & Confidence can be unstable across paraphrases; access and supervision requirements differ. \\
New-knowledge fine-tuning \citep{gekhman-etal-2024-fine} & Controlled closed-book QA experiments vary the amount of unfamiliar factual information in SFT. & The findings motivate knowledge-aware supervision within the studied fine-tuning setup. \\
Known facts \citep{jiang-etal-2024-large} & Examines layer-wise prediction dynamics when questions about the same fact yield correct versus incorrect answers. & Errors can reflect knowledge access rather than absence; measured patterns do not prove universal causal mechanisms. \\
Lost in the Middle \citep{liu-etal-2024-lost} & Evaluates how position of relevant evidence affects long-context use. & Evidence being present is insufficient for reliable use; its tasks are not themselves a comprehensive hallucination taxonomy. \\
Hidden-state limits \citep{servedio-etal-2025-hidden} & Tests factuality encoding on more realistic, model-dependent generated data. & Accuracy on synthetic true/false facts does not establish probe generalization to target-model outputs. \\
Unfamiliar supervision \citep{kang-etal-2025-unfamiliar} & Relates hallucination response patterns to supervision of unfamiliar fine-tuning examples, including reward-model errors. & Factuality alignment depends on how unknowns and reward errors are labeled. \\
Illusion of Progress \citep{janiak-etal-2025-illusion} & Examines lexical-overlap labels and response-length heuristics in detector evaluation. & Semantic validation and simple baselines remain necessary; an LLM judge is not automatically a gold standard. \\
CoT obscures cues \citep{cheng-etal-2025-chain} & Measures how reasoning prompts shift hidden/probability signals and detector effectiveness. & A generation intervention shifts the detector's input distribution, making calibration transfer uncertain. \\
Hallucination tax \citep{song-etal-2025-hallucination} & Studies refusal degradation after reinforcement fine-tuning, primarily on unanswerable mathematics, with reported transfer to factual QA. & Evidence is not a universal result for every factual-generation task. \\
Source preferences \citep{schuster-etal-2026-whose} & Controlled synthetic-source study of source preferences and repetition in conflicting English evidence. & Frequency and source cues influence conflict resolution; synthetic-source preference is not proof of real-world source credibility. \\
\bottomrule\end{tabularx}
\caption{Studies of hallucination mechanisms and reliability in LLMs.}
\label{tab:extra-analysis}
\end{table*}

\begin{table*}[tp]
\centering\small\setlength{\tabcolsep}{5pt}\renewcommand{\arraystretch}{1.18}
\begin{tabularx}{\textwidth}{@{}p{.22\textwidth}p{.34\textwidth}X@{}}\toprule
Study & Mechanism or evaluation design & Comparative implication and boundary \\\midrule
FaithfulRAG \citep{zhang-etal-2025-faithfulrag} & Models fact-level conflicts between retrieved and parametric knowledge before response generation. & Explicit reconciliation complements source preference; source adherence and world correctness remain distinct. \\
CoCoA \citep{khandelwal-etal-2025-cocoa} & Adapts token-level decoding using confidence/context distribution differences. & Adaptive and fixed contrast differ in conflict sensitivity; low-conflict behavior also matters. \\
Context-DPO \citep{bi-etal-2025-context} & Constructs preferences from faithful versus stubborn responses under knowledge conflict. & A context-faithfulness objective differs from truthfulness with respect to external facts. \\
\bottomrule\end{tabularx}
\caption{Additional hallucination mitigation methods for LLMs.}
\label{tab:extra-mitigation}
\end{table*}


\FloatBarrier
\section{A Common Reporting Protocol}\label{app:protocol}
The following protocol is a synthesis of the comparisons in Sections~\ref{sec:detection}--\ref{sec:analysis}, not a benchmark run. It can be used to evaluate whether an intervention in Figure~\ref{fig:mitigation} solves the kind of error illustrated in Figure~\ref{fig:examples} under the target in Figure~\ref{fig:evaluation}.

\begin{table*}[tp]
\centering\small\setlength{\tabcolsep}{4pt}\renewcommand{\arraystretch}{1.14}
\begin{tabularx}{\textwidth}{@{}p{.17\textwidth}p{.24\textwidth}p{.25\textwidth}X@{}}\toprule
Family / example & Required access or supervision & Principal extra resource & Important boundary \\\midrule
Source alignment / AlignScore & Source and output text; trained auxiliary model & Verifier inference over chunks & Auxiliary-model errors and evidence truncation \\
Claim verification / FActScore & Claim extraction, knowledge source, retrieval and verifier & Search and per-claim checks & Retrieval failure differs from disproof \\
Self-checking / SelfCheckGPT & Repeated outputs; no generator internals & Additional generation and comparisons & Repeated falsehood can appear consistent \\
Semantic entropy & Samples and semantic equivalence; likelihoods for weighted estimation & Sampling and clustering & Estimator access differs by variant \\
Probe / Lookback Lens & Attention maps and labeled training spans & Feature extraction and probe inference & Supervised contextual proxy, not proof of truth \\
Factuality alignment / FactAlign & Generator training and sentence-level factual feedback & Verification plus optimization & Reward errors can be reinforced \\
Retrieval control / Self-RAG & Training, retrieval, reflection-token supervision & Training and retrieved candidates & Reflection judgments need auditing \\
Contrastive decoding / CAD & Context-conditioned and context-free logits & Extra distribution at decoding & Greater source adherence can amplify false evidence \\
Layer contrast / DoLa & Intermediate layer predictions & Layer-output computation & Parametric knowledge is not external evidence \\
Representation editing / TruthX & Hidden states and learned auxiliary components & Auxiliary learning and intervention & Transfer beyond training conditions is uncertain \\
Verification / CoVe, CRITIC & Model calls; CRITIC also needs tools & Verification and revision rounds & Isolated prompts are not independent knowledge \\
Selective generation / SEAL & Rejection-token training and regularized decoding & Training; altered answering behavior & Lower risk must be interpreted at matched coverage \\
\bottomrule\end{tabularx}
\caption{Access requirements and computational resources of representative methods. Resource descriptions are qualitative and do not imply measured cost rankings.}
\label{tab:access}
\end{table*}

\subsection{Fix the target before selecting a metric}
Specify whether the target is world correctness, source support, or both; the answer unit; and the treatment of unknown or unverifiable claims. Record the generator/version, prompts, retrieval corpus/date, evidence truncation, decoding parameters, and judge/version. Use the same evidence and question distribution across systems where possible. If a method changes retrieval, separate retrieval quality from generation conditional on the retrieved passages.

For a nonempty extracted claim set $C$ and fixed evidence $E$, a schematic support precision is
\begin{equation}
P_{\mathrm{support}}(C,E)=\frac{\sum_{c\in C}\mathbf{1}[E\models c]}{|C|}.
\end{equation}
This expression describes an aggregation principle, not the complete implementation of every factuality metric. Extraction determines the denominator; duplicate, trivial, subjective, and undecidable claims need explicit handling. An empty answer is not assigned perfect useful factuality by this definition. FActScore and VeriScore motivate claim-level analysis, while Core highlights the role of informative subclaims \citep{min-etal-2023-factscore,song-etal-2024-veriscore,jiang-etal-2025-core}.

\subsection{Separate detection, localization, and generation quality}
For binary detection, report class prevalence, confusion counts, precision/recall/F1, and balanced accuracy, $\tfrac12(\mathrm{TPR}+\mathrm{TNR})$, where appropriate. Fix thresholds on independent validation data. For spans, specify tokenization and exact versus overlap matching. For continuous scores, report discrimination and calibration separately; neither establishes the truth of an individual answer. Inspect false-positive and false-negative examples across error types, evidence lengths, and generator families.

For generation, pair factuality with unique supported content and required-aspect coverage. A source-faithful answer may still omit the answer to the question. Human auditing should sample both high- and low-scoring outputs, use explicit instructions, and report disagreement or adjudication. LongEval provides relevant human-evaluation guidance \citep{krishna-etal-2023-longeval}. Avoid pooling scalar scores whose claim extraction, evidence, or unknown-label policy differs.

\subsection{Measure selection and revision effects}
Let $A_\tau$ be the answers retained by threshold $\tau$ among $N$ questions, and let $e_i$ indicate an erroneous answer under the declared reference frame. For nonempty $A_\tau$, define
\begin{equation}
\begin{aligned}
\mathrm{Coverage}(\tau)&=\frac{|A_\tau|}{N},\\
\mathrm{Risk}(\tau)&=\frac{\sum_{i\in A_\tau}e_i}{|A_\tau|}.
\end{aligned}
\end{equation}
Comparing risk at matched coverage prevents a system from appearing universally better merely by refusing more questions. This response-level selective-evaluation definition is distinct from SEAL's token-level training objective. For revision, count corrected errors, newly introduced errors, retained supported claims, and required information lost. Measure both before and after revision on the same initial drafts.

\subsection{Audit independence, robustness, and cost}
A verifier used for training or candidate selection should not be the sole final judge. Use independently collected labels or an independently audited evaluator. Test natural errors from held-out generators, difficult consistent falsehoods, evidence conflicts, context permutations, paraphrases, and updated facts. Disaggregate performance rather than interpreting an average as proof of transfer. Verify with Caution and temporal benchmark analysis motivate independent error inspection and timestamped evidence \citep{godbole-jia-2025-verify,jiang-etal-2026-benchmarks}.

Table~\ref{tab:access} summarizes method access and deployment trade-offs. Report offline annotation and training costs separately from inference calls, sampled tokens, retrieval requests, tool invocations, latency, and memory. A black-box method can be expensive, while an internal method can be unavailable even if computationally economical. Do not infer a fixed latency multiplier from the number of model distributions or stages without measurement. Comparisons should either match budgets or present a quality--cost trade-off.

% Genuine citations printed in imported figures; registered for BibTeX.
\nocite{adlakha-etal-2024-evaluating,azaria-mitchell-2023-internal,bang-etal-2025-hallulens,bao-etal-2025-faithbench,cao-etal-2022-hallucinated,cheang-etal-2026-llms,chuang-etal-2024-lookback,cohen-etal-2023-lm,dhuliawala-etal-2024-chain,dziri-etal-2022-evaluating,dziri-etal-2022-faithdial,extcritic2024,extdola2024,extselfrag2024,extsemanticentropy2024,extsemanticuncertainty2023,fabbri-etal-2022-qafacteval,fan-etal-2026-refl,fatahi-bayat-etal-2025-factbench,feng-etal-2024-dont,gema-etal-2025-decore,godbole-jia-2025-verify,harary-etal-2026-prefixnli,huang-chen-2024-factalign,huang-etal-2025-alleviating,huang-etal-2025-improving,ji-etal-2024-anah,jiang-etal-2023-active,jiang-etal-2026-benchmarks,li-etal-2023-halueval,li-etal-2024-dawn,lin-etal-2022-truthfulqa,liu-etal-2024-universal,malaviya-etal-2024-expertqa,manakul-etal-2023-selfcheckgpt,min-etal-2023-factscore,niu-etal-2024-ragtruth,ravichander-etal-2025-halogen,samarinas-etal-2025-beyond,semnani-etal-2023-wikichat,shi-etal-2024-trusting,song-etal-2024-veriscore,tang-etal-2024-tofueval,wang-etal-2024-factuality,wang-etal-2025-astute,wu-etal-2024-synchronous,xue-etal-2025-ualign,zha-etal-2023-alignscore,zhang-etal-2023-enhancing-uncertainty,zhang-etal-2023-sac3,zhang-etal-2024-self,zhang-etal-2024-truthx,zhang-etal-2026-stable,zhao-etal-2025-response,zhou-etal-2023-context}

